\section{Reproducing the examples}
\label{app:repro}

All listings in this paper are drawn from the repository's doctested
tutorials, examples, and tests of \pqlang{} 0.8.0.  The commands below
regenerate the artifacts referenced in the text:

\begin{lstlisting}[style=pyqbase, language=bash, caption={Reproduction commands (generation and export have no third-party requirements; native cross-checks require the pinned backend environments).}, label={lst:repro}]
# tutorials and core semantics
python -m unittest discover -s tests/core -v

# gallery: 22 cases, RIR + OriginIR-ext + strict gate set exports
python examples/algorithm_gallery.py            # add --native for backend agreement

# oracle/protocol/QODE/QHAM demonstration programs
python examples/algorithm_contracts.py
python examples/ode_input_models.py             # 11 input-model cases
python examples/input_models.py                 # 14 input-model variants over 4 families
python examples/general_qham.py                 # forced Burgers at order 2
python examples/math_functions.py

# static checks
ruff check src tests examples

# numerical verification campaign (Section 9.3): 473 experiments in 13
# groups, writes machine-readable reports to out/verification/
PYTHONPATH=src <backend-python> tools/run_verification.py
\end{lstlisting}

The native configuration runs the same examples under interpreters
providing \texttt{pysparq} and \texttt{uniqc}; revision pins are recorded in
the repository's backend revision ledger.  The verification campaign
requires both backends: the PySparQ path executes RIR natively through
\texttt{pysparq.rir}, and the OriginIR-ext path uses UnifiedQuantum's
\texttt{Circuit.to\_matrix()}---both contributed as part of this work.

\section{The four solver families in code}
\label{app:families}

This appendix shows, for each of the four solver territories of
\secref{sec:ode}, a complete input program drawn verbatim from the
repository's doctested examples (comments translated from Chinese;
artifact bookkeeping is elided at the marked points).  The point of
interest is not the numerics but what stays \emph{open} and what is
\emph{swappable} in each family: every listing contains at least one
\texttt{abstract\_*} declaration that the same program could bind
differently, and every solver entry accepts its kernel configuration as an
argument (R1, R2, R5).

\subsection{QHAM: from a symbolic PDE to a solver-ready input model}
\label{app:qham}

\lstref{lst:qham} is the repository's forced-Burgers example
(\texttt{examples/general\_qham.py}) in full.  The equation is entered
\emph{symbolically}---a \texttt{PolynomialPDE} over a \texttt{Field}
expression, with the forcing term declared as named data
(\texttt{Known("forcing")}) rather than a baked-in array.  The homotopy
plan fixes the linearization order ($m=2$) and the homotopy parameter
($\eta=-0.4$); the structured finite-difference bindings carry the
stencil-structured ports and the lifted initial state.  The result,
\texttt{qode\_input}, is an open QODE input model: it can be handed to any
solver family, and the physically owned pieces (coefficient block-encodings,
forcing data) remain bindable separately.

\begin{lstlisting}[style=pyqpython, caption={QHAM end to end: symbolic polynomial PDE, homotopy-analysis plan, structured discretization, and a solver handed in from outside (verbatim from the doctested example \texttt{examples/general\_qham.py}; comments translated).}, label={lst:qham}]
from functools import partial

from pyqecclang.algorithms.hamiltonian import taylor_hamiltonian
from pyqecclang.algorithms.ode import linear_qode
from pyqecclang.applications.qham import (
    Discretization,
    Field,
    Grid,
    Known,
    PolynomialPDE,
    QHAMPlan,
    qham_input_model,
    structured_fd_bindings,
)

u = Field("u")
pde = PolynomialPDE.from_equations(
    {"u": 0.1 * u.d("x", 2) - u * u.d("x") + Known("forcing")},
    label="forced_burgers",
)
plan = QHAMPlan(pde, order=2)
grid = Grid(("x",), (4,), (1.0,), boundary="periodic")
space = Discretization(pde, grid, {"forcing": [0.05, 0, -0.05, 0]})
initial = space.encode_fields({"u": [0.1, 0.2, 0, -0.1]})
bindings = structured_fd_bindings(space, initial)
qode_input = qham_input_model(plan, bindings, eta=-0.4)

# Schrodingerization accepts a general generator; recovery-interval and
# time-approximation validation remains a pending item.
solver = linear_qode(
    "schrodingerization", hamiltonian_function=partial(taylor_hamiltonian, degree=1)
)
solution = qode_input.solve(solver, 0.01)
\end{lstlisting}

\subsection{Carleman: nonlinear lift with a swappable linear solver}
\label{app:carleman}

\lstref{lst:carleman} (case~6 of \texttt{examples/ode\_input\_models.py})
runs the inviscid Burgers equation through the Carleman pipeline.  Each
polynomial coefficient $F_p$ of the lifted generator is declared as an
\emph{open} block-encoding slot (\texttt{BurgersF$p$}) and bound to the
operation produced by the stencil ports, so the discretization could be
replaced without touching the problem assembly.  The lifted problem is a
\texttt{PolynomialODE}; the solver,
\texttt{carleman\_qode(linear\_solver=..., cutoff=2)}, receives the linear
solver as an argument, and the loop solves the \emph{same} nonlinear
problem twice---once through Schr\"odingerization and once through
norm-shifted LCHS---which is exactly the ``Carleman lands in the same
linear-solver position'' claim of \secref{sec:ode-solvers}.

\begin{lstlisting}[style=pyqpython, caption={Carleman lift of Burgers: open coefficient block-encodings, an open lifted initial state, and a linear solver injected per call (verbatim from case~6 of \texttt{examples/ode\_input\_models.py}; comments translated, artifact serialization elided).}, label={lst:carleman}]
u = Field("u")
pde = PolynomialPDE.from_equations({"u": 0.1 * u.d("x", 2) - u * u.d("x")})
space = Discretization(pde, grid)
ports = structured_fd_bindings(space, space.encode_fields({"u": [0.1, 0.2, 0, -0.1]}))
concrete = polynomial_from_bindings(ports)
coefficient_slots, bindings = [], {}
for order, coefficient in concrete.coefficients:
    name = "BurgersF" + str(order)
    slot = abstract_block_encoding(
        name, coefficient.width, coefficient.signal_qubits, coefficient.alpha
    )
    coefficient_slots.append((order, slot))
    bindings[name] = coefficient.operation
initial_slot = abstract_state_prep(
    "BurgersInitial", concrete.width, concrete.initial.work_width
)
bindings["BurgersInitial"] = concrete.initial.operation
problem = PolynomialODE(
    concrete.width, tuple(coefficient_slots), initial_slot, concrete.initial_norm
)
recovery = {}
for method, solver in (
    ("schrodingerization", schrodinger),
    ("shifted_lchs", shifted_solver(lchs, recovery)),
):
    nonlinear_solver = partial(carleman_qode, linear_solver=solver, cutoff=2)
    state = qpde_solver(nonlinear_solver)(PDEInput(problem, "burgers"), time)
\end{lstlisting}

\subsection{LCHS: one open program, two data-path realizations}
\label{app:lchs}

\lstref{lst:lchs} condenses three cases of
\texttt{examples/ode\_input\_models.py}.  The protocol object is built by
the uniform \texttt{linear\_qode} front with two swappable arguments: the
quadrature plan (Cauchy, cutoff 1) and the injectable Hamiltonian-function
kernel.  The first problem declares only an abstract generator block-encoding;
the third declares a rotation-angle database, and the \emph{same} open
program is then closed twice---once against a gate realization of the
database and once against a QRAM database with a runtime memory
table---with no change to the program itself (R2).

\begin{lstlisting}[style=pyqpython, caption={LCHS: the protocol takes its quadrature plan and kernel as arguments; one open program is bound to a gate database or to a QRAM database with a runtime table (verbatim from \texttt{examples/ode\_input\_models.py}, inside \texttt{main()}; comments translated, artifact records elided where marked).}, label={lst:lchs}]
time = 0.05
hamiltonian_function = partial(taylor_hamiltonian, degree=1)
quadrature = QuadraturePlan.cauchy(cutoff=1, spacing=1.0)
lchs = linear_qode("lchs", plan=quadrature, hamiltonian_function=hamiltonian_function)

# Generator block-encoding given directly; G = -I.  The work qubits are
# reserved for the QRAM initial-state realization used below.
initial = abstract_state_prep("Initial", 1, work_width=3)
initial_gate = gate_state_prep([1, 1], work_width=3)
generator = abstract_block_encoding("Generator", 1, 0, 1.0)
state = lchs(generator, initial, time)
# ... artifact record elided ...

# Same open graph, bound once against a gate database and once against
# a QRAM database plus state preparation.  The data are rotation-angle
# integer encodings; A = diag(1, cos(pi/4)), G = -A.
angles = abstract_database("DiagonalAngles", 1, 2)
a = diagonal_block_encoding(angles, alpha=1.0)
state = lchs(scale(-1, a), initial, time)
records.append(
    write(
        "given_xor_gate_lchs",
        state,
        {
            "DiagonalAngles": gate_database(1, 2, [0, 1]).operation,
            "Initial": initial_gate.operation,
        },
    )
)
# ... and the QRAM realization of the same slots:
records.append(
    write(
        "given_xor_qram_lchs",
        state,
        {
            "DiagonalAngles": Binding(
                qram_database(1, 2).operation, {"table": "diagonal_angles"}
            ),
            "Initial": Binding(qram_state_prep(1, 2).operation, {"angles": "initial_angles"}),
        },
        memory={
            "diagonal_angles": [0, 1],
            "initial_angles": qram_state_angles([1, 1], 2),
        },
    )
)
\end{lstlisting}

\subsection{CBMD: the contour family under a dissipative shift}
\label{app:cbmd}

The CBMD family is reached through the same uniform front---the only
difference from \lstref{lst:qode-swap} is the family name string.
Because the contour construction requires a strictly dissipative
generator, the QHAM input model of \lstref{lst:qham} first applies
its \texttt{dissipative\_shift} adaptation (the explicit global shift
$G \mapsto G - \mu I$ with the host tracking the norm bookkeeping); the
shifted model then solves against \texttt{"cbmd"} with the default contour
configuration and the same injectable kernel.  \lstref{lst:cbmd}
is the final block of \texttt{examples/general\_qham.py}: two lines
separate the Schr\"odingerization solution of \lstref{lst:qham}
from the CBMD solution of the same problem.

\begin{lstlisting}[style=pyqpython, caption={CBMD: the same QHAM input model after \texttt{dissipative\_shift}, solved by naming the family (verbatim from the final block of \texttt{examples/general\_qham.py}; comments translated).}, label={lst:cbmd}]
# For CBMD/LCHS, apply the explicit global shift that makes the
# generator dissipative first.
shifted = qode_input.dissipative_shift()
cbmd_solution = shifted.solve(
    linear_qode("cbmd", hamiltonian_function=partial(taylor_hamiltonian, degree=1)), 0.01
)
\end{lstlisting}

Across the four families the pattern is constant: symbolic or structured
problem entry, oracle slots declared open at the point where physics or
hardware ownership is contested, and solver families selected by name with
their approximation kernels injected as ordinary arguments.  None of the
listings constructs a circuit; the closest any of them comes to circuit
vocabulary is the binding dictionary that names which realization closes
which slot.

\subsection{One problem, many data paths: the input-model matrix}
\label{app:input-models}

Sections~\ref{app:qham}--\ref{app:cbmd} fixed one input model per family;
\texttt{examples/input\_models.py} instead holds the \emph{problem} fixed
and varies the data path, producing the fourteen artifacts of
\tabref{tab:input-models}.  Three input-model families
appear.  \emph{Structured gate stencils} bake the stencil and coefficient
data into gate logic (the baseline of the previous subsections).
\emph{Spectral embeddings} replace the structured construction by an
exact Pauli-LCU decomposition of each port matrix
(\texttt{gate\_bindings}) or of a given generator
(\texttt{matrix\_pauli\_encoding}); the same decomposition with a
\texttt{qram\_prepare} table turns the Pauli coefficients into QRAM
runtime data.  \emph{QRAM data paths} leave the data in named resources:
coefficient values become angle words in an open database slot
(\texttt{qram\_coefficient\_encoding}, plugged into the
\texttt{coefficient\_encoder} hook of the structured bindings), the
grid's sparse-access position and entry oracles become QRAM tables, and
the initial state is loaded through \texttt{qram\_state\_prep}.

\begin{table}[t]
  \centering
  \caption{The input-model matrix of \texttt{examples/input\_models.py}.
  All artifacts share the same small instances (forced Burgers for QHAM,
  Burgers with a space-dependent viscosity for the Carleman QRAM row, a
  given $2\times2$ generator, and a four-node ring Laplacian).  ``Open
  slots'' lists what remains unbound in the exported program; ``runtime
  tables'' lists the QRAM resources present after closing (memory
  supplied separately).  Gate rows are exact; angle-table rows carry the
  declared \texttt{angle\_width} quantization.}
  \label{tab:input-models}
  \footnotesize
  \setlength{\tabcolsep}{3.5pt}\begin{tabular}{@{}p{1.5cm}p{4.3cm}p{3.9cm}p{4.3cm}@{}}
    \toprule
    \tabhead{Family} & \tabhead{Input model} & \tabhead{Open slots} & \tabhead{Runtime tables} \\
    \midrule
    QHAM & structured stencil ports (gate) & --- & --- \\
    QHAM & spectral ports (Pauli LCU) & --- & --- \\
    QHAM & coefficient-angle db + declared initial & coefficient db, initial & \texttt{coeff\_angles}, \texttt{initial\_angles} (gate table or QRAM) \\
    \midrule
    Carleman & structured stencil ports & $F_1$, $F_2$, initial & --- \\
    Carleman & spectral ports (Pauli LCU) & $F_1$, $F_2$, initial & --- \\
    Carleman & viscosity field via angle database + QRAM initial & $\nu$ db, $F_1$, $F_2$, initial & \texttt{nu\_angles}, \texttt{initial\_angles} \\
    \midrule
    LCHS & given spectral generator (gate) & Generator, Initial & --- \\
    LCHS & given spectral generator, QRAM PREPARE & Generator, Initial & \texttt{spectral\_angles} \\
    LCHS & QRAM grid stencil + QRAM initial & position, entry, initial & \texttt{positions}, \texttt{inverse\_positions}, \texttt{entries}, \texttt{grid\_angles} \\
    \midrule
    CBMD & given spectral generator (gate) & Generator, Initial & --- \\
    CBMD & given spectral generator, QRAM PREPARE & Generator, Initial & \texttt{spectral\_angles} \\
    CBMD & QRAM grid stencil + QRAM initial & position, entry, initial & \texttt{positions}, \texttt{inverse\_positions}, \texttt{entries}, \texttt{grid\_angles} \\
    CBMD & QHAM coefficient input after \texttt{dissipative\_shift} & coefficient db, initial & \texttt{coeff\_angles}, \texttt{initial\_angles} \\
    \bottomrule
  \end{tabular}
\end{table}

\lstref{lst:input-qham} shows the QHAM row: the coefficient database and
the initial state are declared open, the program is solved once, and the
\emph{same} open program closes against a gate truth table or against a
QRAM database with a runtime table.  \lstref{lst:input-grid} shows the
grid row: the four-node ring Laplacian is exposed as a sparse-access
oracle pair whose position and entry tables are QRAM resources, wrapped
as a \texttt{DiscretePDE} and solved unchanged by two different families.

\begin{lstlisting}[style=pyqpython, caption={QHAM with a coefficient-angle database: one open program, closed against a gate database or a QRAM database (condensed from \texttt{examples/input\_models.py}; comments translated, slot discovery and artifact records elided).}, label={lst:input-qham}]
profile = [0.1, 0.2, 0.15, 0.05]  # QRAM angle tables take non-negative amplitudes
qinitial = space.encode_fields({"u": profile})
encoder = partial(qram_coefficient_encoding, angle_width=8)
open_bindings = structured_fd_bindings(space, qinitial, coefficient_encoder=encoder)
open_bindings = QHAMBindings(
    open_bindings.state_width,
    open_bindings.ports,
    abstract_state_prep("QhamInitial", open_bindings.state_width, 10),
    open_bindings.initial_norm,
)
model = qham_input_model(plan, open_bindings, eta=-0.4)
state = model.solve(schrodinger, 0.01)
# ... the coefficient database slot is discovered from the open program ...
records.append(write("qham_coeff_gate", state, {
    db_slot: gate_database(2, 8, word_table).operation,
    "QhamInitial": gate_state_prep(profile, work_width=10).operation,
}))
records.append(write("qham_coeff_qram", state, {
    db_slot: Binding(qram_database(2, 8).operation, {"table": "coeff_angles"}),
    "QhamInitial": Binding(qram_state_prep(2, 8).operation, {"angles": "initial_angles"}),
}, memory={
    "coeff_angles": word_table,
    "initial_angles": qram_state_angles(profile, 8),
}))
\end{lstlisting}

\begin{lstlisting}[style=pyqpython, caption={A QRAM-served grid discretization: the ring Laplacian as a sparse-access oracle pair with QRAM position and entry tables, plus a QRAM initial state; \texttt{make\_qpde} threads it into two solver families unchanged (verbatim from \texttt{examples/input\_models.py}; comments translated, runtime tables elided).}, label={lst:input-grid}]
fmt = FixedFormat(4, 1)
access = abstract_sparse_access("RingStencil", 2, fmt.width, sparsity=3)
laplacian = real_symmetric_sparse_encoding(access, fmt, 2.0, diagonal_nonnegative=True)
heat = DiscretePDE(
    scale(-0.1, laplacian),
    abstract_state_prep("GridInitial", 2, 10),
    label="ring_heat",
)
# ... forward/inverse position tables, entry table, angle table elided ...
for family, solver in (("lchs", lchs), ("cbmd", cbmd)):
    state = make_qpde(solver)(heat, time)
\end{lstlisting}

The matrix makes the language-level point of
\secref{sec:input-models} concrete: input models are a property
of the \emph{program}, not of the solver.  Nothing in the LCHS, CBMD,
Schr\"odingerization, or Carleman entries changes as the data path moves
from gate tables to QRAM resources; the difference is entirely carried by
which slots stay open, which resources the closed program declares, and
which runtime tables the host must supply.  The limitations of the
angle-table realizations are recorded honestly in the artifacts: QRAM
state tables accept non-negative real amplitudes, and angles are
quantized to \texttt{angle\_width} bits.  Note that the gate and QRAM
closures of the same program encode \emph{identical} data---the gate
truth table stores the very angle words the QRAM table would---so the two
realizations agree exactly, and the declared quantization is the only
approximation relative to the exact coefficients.

\section{Glossary of recurring identifiers}
\label{app:glossary}

The code surface on which this paper relies is deliberately small---fifteen
names cover every program printed here, from the minimal Bell listing to
the QFVM pipeline.  Two conventions organize them.  First, the layering of
\figref{fig:stack} is visible in the prefixes: \texttt{Builder} and
\texttt{Operation} belong to the construction layer, unadorned names such
as \texttt{Module}, \texttt{Ref}, and \texttt{Load} are RIR concepts, and
the \texttt{abstract\_*} family lives in the algorithm libraries, which
declare what a solver \emph{needs} without constructing anything.
Second, solver entry points are uniform fronts rather than class
hierarchies: \texttt{linear\_qode} and \texttt{carleman\_qode} take the
family as data (a name string) together with injectable protocol
arguments, which is what allows the one-line family swaps shown in
\lstref{lst:qode-swap} and Appendix~\ref{app:families}.

\begin{table}[ht]
  \centering
  \caption{Recurring identifiers used throughout the paper.}
  \label{tab:glossary}
  \small
  \setlength{\tabcolsep}{4pt}\setlength{\tabcolsep}{4pt}\begin{tabular}{@{}p{3.3cm}p{10.4cm}@{}}
    \toprule
    \tabhead{Identifier} & \tabhead{Meaning} \\
    \midrule
    \texttt{Builder} & embedded module constructor over typed register views \\
    \texttt{Operation} & finished builder product; satisfies unitary / state-preparation / block-encoding protocols \\
    \texttt{Module} (RIR) & named typed interface + body (or absent body $=$ open oracle slot) \\
    \texttt{Ref} / \texttt{Span} & declarative register view; $(\text{register}, \text{start}, \text{width})$ \\
    \texttt{QRAM} & typed resource: address width, data width, runtime table binding \\
    \texttt{Load} & reversible memory access $|a\rangle|d\rangle \mapsto |a\rangle|d \oplus M[a]\rangle$ \\
    \texttt{bind} / \texttt{unresolved} & oracle linking pass / open-slot gap analysis \\
    \texttt{BlockEncoding} & operation + normalization $\alpha$ + capabilities \\
    \texttt{QODEProblem} & linear ODE input object (generator BE, initial state, promises) \\
    \texttt{linear\_qode} & uniform entry over LCHS / Schr\"odingerization / CBMD with injectable Hamiltonian-function kernel \\
    \texttt{carleman\_qode} & nonlinear path: polynomial ODE $\rightarrow$ cutoff-$K$ lift $\rightarrow$ any linear solver \\
    \texttt{qram\_coefficient\_}\\ \texttt{encoding} & coefficient diagonal as an open angle database, bound to a gate or QRAM realization \\
    \texttt{make\_costa\_qlss} /\\ \texttt{make\_cks\_qlss} & the two QLSS kernels (interchangeable) \\
    \texttt{SpectralPromise} & norm / smallest-singular-value bounds with provenance \\
    \texttt{compile\_function} & pure-Python function $\rightarrow$ SSA MIR $\rightarrow$ reversible RIR module \\
    \texttt{quantikz} & circuit-diagram export; calls as named boxes, open slots dashed \\
    \texttt{export\_originir} & modular textual export (module calls preserved) \\
    \bottomrule
  \end{tabular}
\end{table}

The table doubles as a stability statement: these names constitute the
public surface of \pqlang{} 0.8.0 on which every doctest, gallery case,
and catalog entry of Appendix~\ref{app:repro} is pinned.  Readers who want
the next level of detail---constructor signatures, protocol definitions,
and the per-family configuration records---will find them in the API
reference shipped with the distribution, which is generated from the same
docstrings the examples exercise.
